\documentclass\setlength{\headheight}{14.49998pt}

% Packages
\usepackage[utf8]{inputenc}
\usepackage[T1]{fontenc}
\usepackage{mathptmx}
\usepackage{graphicx}
\usepackage{amsmath,amssymb,amsfonts}
\usepackage{booktabs}
\usepackage{longtable}
\usepackage{listings}
\usepackage{xcolor}
\usepackage{hyperref}
\usepackage{geometry}
\usepackage{fancyhdr}
\usepackage{setspace}
\usepackage{etoolbox}

% Page setup
\geometry{margin=1.2in}
\onehalfspacing

% Colors for code
\definecolor{codegreen}{rgb}{0,0.6,0}
\definecolor{codegray}{rgb}{0.5,0.5,0.5}
\definecolor{codepurple}{rgb}{0.58,0,0.82}
\definecolor{backcolour}{rgb}{0.95,0.95,0.92}

% Code listing style
\lstdefinestyle{pythonstyle}{
    backgroundcolor=\color{backcolour},
    commentstyle=\color{codegreen},
    keywordstyle=\color{blue},
    numberstyle=\tiny\color{codegray},
    stringstyle=\color{codepurple},
    basicstyle=\ttfamily\small,
    breaklines=true,
    frame=single,
    numbers=left,
    numbersep=5pt
}

% Header/Footer
\pagestyle{fancy}
\fancyhf{}
\rhead{\thepage}
\lhead{MASSIF Research Report --- DUBITO Inc. / Ergo Sum}
\chead{}

% Title definition
\title{\textbf{Research Report on Hidden-State Geometry and the Building of a Rigorous Empirical Infrastructure for Latent Trajectory Analysis Falsifying the Early-Warning Hypothesis}\\[1em]
\large DUBITO Inc. / Ergo Sum AGI Safety Systems\\
\small Research Memorandum -- solis@dubito-ergo.com}

\author{Daniel Solis}

\date{11. May 2026}

\begin{document}

\maketitle

\begin{center}
\vspace*{2em}
\textbf{Abstract}
\end{center}

\noindent
This report documents the full methodological trajectory of the MASSIF research program, beginning with speculative CQFT-inspired geometric theories and evolving into a controlled empirical investigation of transformer hidden-state dynamics. The original hypothesis proposed that unsafe latent cognition might produce measurable geometric precursors before overt unsafe semantics appear in generated text.

Initial exploratory experiments appeared to support this claim, suggesting that hidden-state geometry diverged prior to explicit unsafe semantic manifestation. However, subsequent controlled experiments, expanded datasets, negative controls, and revised semantic detection methodologies failed to reproduce the effect. The apparent precursor signal was determined to be a methodological artifact arising primarily from delayed semantic labeling rather than genuine predictive latent geometry.

Although the central precursor hypothesis was not supported, the research program produced several meaningful findings. Transformer hidden states do exhibit robust geometric divergence following semantic branching, larger models display qualitatively stronger divergence structures, and the developed framework provides a viable infrastructure for large-scale latent trajectory telemetry research.

The report argues that the scientific value of the project lies not in confirmation of the original hypothesis, but in the disciplined falsification process that transformed speculative geometric intuition into a substantially more rigorous empirical framework.

\medskip
\hrule
\medskip

\section{1. Origins: CQFT and the Search for Geometric Criticality}

The MASSIF trajectory originated within the broader CQFT/PFTC research program, whose initial objective was to investigate whether coherent field geometries could induce substrate-selective criticality and self-organizing behavior.

The early conceptual chain was approximately:

\begin{center}
\begin{minipage}{0.6\textwidth}
\begin{align*}
\text{C} &\rightarrow \text{G} \\
\text{Geometry} &\rightarrow \text{Stability} \\
\text{Critical Topology} &\rightarrow \text{Coherent Emergence} \\
\text{Phase Transitions} &\rightarrow \text{Substrate Intelligence}
\end{align*}
\end{minipage}
\end{center}

This framework drew inspiration from dissipative systems, topological criticality, geometric coherence, and speculative substrate-selection principles.

The original experiments attempted to operationalize these ideas through fractal and topological simulations intended to reveal phase-transition behavior analogous to Penrose-style geometric criticality.

\textbf{The result was negative.}

No robust phase transition emerged. No $\phi$-locking appeared. No topology-sensitive amplification survived controlled testing. The original ``criticality ladder'' interpretation failed under operationalization.

In retrospect, this failure became the most scientifically productive event in the program.

\medskip
\hrule
\medskip

\section{2. The Transformer Pivot}

Following the collapse of the original criticality hypothesis, attention shifted toward transformer hidden-state analysis.

The conceptual translation was cautious and incomplete:

\begin{center}
\begin{tabular}{ll}
Hidden-state manifold & $\approx$ field structure \\
Attention geometry & $\approx$ latent topology \\
Trajectory evolution & $\approx$ dynamical geometry
\end{tabular}
\end{center}

This was never interpreted as literal confirmation of CQFT. Rather, CQFT became a heuristic inspiration for investigating whether transformer latent spaces possessed measurable geometric organization.

This distinction proved crucial. The project ceased to be an attempt to validate speculative consciousness theory and instead evolved into a program of empirical latent telemetry research.

\medskip
\hrule
\medskip

\section{3. Early Positive Findings}

Initial experiments examined hidden-state structure across multiple GPT-2 variants.

Unexpectedly, coherent generations displayed substantially greater geometric organization than randomized controls. Geometry metrics appeared to scale with model size:

\begin{equation}
0.6004 \rightarrow 0.7588 \rightarrow 1.1992 \rightarrow 1.4209
\end{equation}

This contradicted the original expectation that coherence would flatten geometric structure through entropy minimization.

Subsequent experiments introduced behavioral categories: truthful, deceptive, sycophantic, hallucinatory, and jailbreak-oriented. Hidden-state geometry appeared to vary systematically between these behavioral regimes.

At the time, this suggested the possibility that unsafe cognition might possess measurable latent geometric signatures.

\medskip
\hrule
\medskip

\section{4. The Emergence of the Early-Warning Hypothesis}

The research focus gradually shifted toward a more operational question:

\begin{quote}
\textbf{Research Question:} Can hidden-state geometry detect unsafe trajectories before unsafe text appears?
\end{quote}

This became formalized through the proposed metric:

\begin{equation}
\text{EPH} = t_{\text{semantic}} - t_{\text{geometric}}
\end{equation}

where:
\begin{itemize}
    \item $t_{\text{geometric}}$ represents the first detectable geometric divergence,
    \item $t_{\text{semantic}}$ represents the first overt unsafe semantic manifestation.
\end{itemize}

Early exploratory experiments appeared promising. In paired prompt tests, hidden-state divergence sometimes appeared several tokens before explicit unsafe phrases emerged.

This produced the tentative hypothesis that unsafe trajectories might possess latent geometric precursors.

At this stage, however, the experiments remained exploratory:
\begin{itemize}
    \item very small sample sizes,
    \item no negative controls,
    \item unstable PCA regimes,
    \item ad hoc composite metrics,
    \item and weak semantic alignment procedures.
\end{itemize}

These limitations would later become decisive.

\medskip
\hrule
\medskip

\section{5. MASSIF and the Move Toward Statistical Rigor}

The need for systematic scaling led to the design of MASSIF:

\subsection*{Massive Automated Statistical Sweep Framework}

The goal was to industrialize latent trajectory analysis through:

\begin{center}
Prompt Generation $\rightarrow$ Counterfactual Branching $\rightarrow$ Hidden-State Extraction $\rightarrow$ Geometry Telemetry $\rightarrow$ Statistical Evaluation $\rightarrow$ Predictive Testing
\end{center}

At this stage the project underwent an important philosophical transformation.

The research ceased to pursue:
\begin{itemize}
    \item consciousness validation,
    \item quantum cognition,
    \item Penrose equivalence,
    \item or metaphysical interpretations.
\end{itemize}

Instead, it reframed itself as:

\begin{quote}
\textbf{Geometric Runtime Telemetry for Transformer Latent Dynamics}
\end{quote}

This reframing dramatically improved the program's methodological clarity.

\medskip
\hrule
\medskip

\section{6. The Controlled Experiments}

MASSIF v0.3 introduced several major corrections absent from earlier notebooks:
\begin{itemize}
    \item larger sample sizes,
    \item negative controls,
    \item lexical semantic detection,
    \item bootstrap confidence intervals,
    \item permutation testing,
    \item metric cross-application,
    \item and control-condition comparisons.
\end{itemize}

Most importantly, semantic onset was redefined.

Earlier notebooks used keyword heuristics such as \texttt{"exploit"}, \texttt{"manipulate"}, \texttt{"deceive"} to determine semantic divergence. This introduced a systematic delay between lexical branching and semantic detection.

The controlled experiments instead measured divergence beginning from the first lexical branch itself.

This methodological change fundamentally altered the results.

\medskip
\hrule
\medskip

\section{7. Collapse of the Precursor Hypothesis}

Under controlled conditions, the original early-warning hypothesis failed.

The measured quantity:

\begin{equation}
\Delta t = t_{\text{semantic}} - t_{\text{geometric}}
\end{equation}

was consistently negative across conditions.

Safe-versus-unsafe comparisons produced:

\begin{equation}
\text{Mean } \Delta t \approx -2.45 \text{ tokens}
\end{equation}

indicating that geometry diverged \emph{after} lexical/semantic branching rather than before it.

Crucially, negative controls exhibited similar behavior:

\begin{table}[h]
\centering
\caption{Early Warning Detection Results by Condition}
\label{tab:early_warning}
\begin{tabular}{@{}lcc@{}}
\toprule
Condition & Mean $\Delta t$ & Interpretation \\
\midrule
Safe vs Unsafe & $-2.45$ & Geometry trails semantics \\
Safe vs Safe & $-6.50$ & Null condition (baseline) \\
Unsafe vs Unsafe & $-2.50$ & Null condition (baseline) \\
\bottomrule
\end{tabular}
\end{table}

This result was devastating for the alignment-specific interpretation.

If unsafe geometry were uniquely predictive, safe-versus-unsafe trajectories should have separated clearly from ordinary semantic branching controls.

They did not.

The conclusion became unavoidable:

\begin{quote}
\textbf{The original precursor finding was primarily an artifact of semantic detection methodology.}
\end{quote}

\medskip
\hrule
\medskip

\section{8. The G-Index Reassessment}

An additional revelation emerged through cross-application testing.

The original composite geometry metric:

\begin{lstlisting}[language=Python, style=pythonstyle, caption={Original G-Index Formula}, label=lst:gindex]
G = explained * 10 + anisotropy + geom_entropy
    + spread * 0.05
    - smoothness * 0.03
    - curvature * 2
\end{lstlisting}

had initially appeared to capture meaningful geometric structure.

However, when applied under controlled conditions, the G-index shifted timing estimates by only approximately $+0.33$ tokens relative to simple cosine divergence.

This strongly suggested that most detectable signal resided in ordinary hidden-state separation rather than sophisticated manifold geometry.

The implication was profound: the elaborate geometric composite added little predictive information beyond basic representational divergence.

\medskip
\hrule
\medskip

\section{9. What Survived Falsification}

Although the central precursor hypothesis failed, several important findings remained intact.

\subsection{9.1 Geometric Divergence Is Real}

Transformer hidden states do reorganize geometrically following semantic branching. The divergence is measurable, robust, and reproducible.

\subsection{9.2 Model Scale Matters}

Larger models exhibited qualitatively stronger divergence behavior, including:
\begin{itemize}
    \item sharper trajectory separation,
    \item stronger anisotropy,
    \item and near-antipodal representation states.
\end{itemize}

This suggests scaling-related changes in latent manifold organization.

\subsection{9.3 Geometry Appears Reactive Rather Than Predictive}

The evidence now supports the following interpretation:

\begin{center}
\begin{minipage}{0.7\textwidth}
Semantic Branching $\xrightarrow{\text{}}$ Latent Geometric Amplification
\end{minipage}
\end{center}

Rather than:

\begin{center}
\begin{minipage}{0.7\textwidth}
Latent Geometry $\xrightarrow{\text{}}$ Semantic Prediction
\end{minipage}
\end{center}

This distinction fundamentally changes the interpretation of the results.

\medskip
\hrule
\medskip

\section{10. Methodological Lessons}

The most important scientific contribution of MASSIF may ultimately be methodological rather than empirical.

The project demonstrated several critical principles for latent-space research:

\subsection{10.1 Semantic Operationalization Dominates Outcomes}

Small changes in semantic labeling strategy can completely reverse temporal conclusions.

\subsection{10.2 Negative Controls Are Essential}

Without safe-safe and unsafe-unsafe controls, ordinary branching artifacts can masquerade as alignment-specific effects.

\subsection{10.3 Exploratory Geometry Is Highly Vulnerable to Confounds}

Sequence length, lexical divergence, entropy variation, and PCA instability can all produce misleading signals.

\subsection{10.4 Falsification Is Productive}

The controlled experiments substantially improved the scientific integrity of the framework by eliminating unsupported interpretations.

\medskip
\hrule
\medskip

\section{11. Current Interpretation}

The current evidence does \textbf{not} support the claim that unsafe latent geometry predictively precedes unsafe semantic manifestation.

However, the work does support the following revised interpretation:

\begin{quote}
\textbf{Transformer hidden-state trajectories undergo measurable geometric reorganization following semantic divergence, and larger models exhibit increasingly strong and structured divergence behavior.}
\end{quote}

This is a narrower claim than originally proposed, but a substantially more defensible one.

\medskip
\hrule
\medskip

\section{12. Future Directions}

Several research directions remain scientifically viable:

\begin{itemize}
    \item scaling laws of latent divergence,
    \item geometry of semantic stabilization,
    \item trajectory amplification dynamics,
    \item latent attractor formation,
    \item post-branch representation polarization,
    \item and comparative topology across aligned versus uncensored models.
\end{itemize}

The original dream of geometric early-warning telemetry was not validated.

But the research process itself revealed something perhaps equally important:

\begin{quote}
\textbf{That interpretability research must survive controlled falsification before its claims can be trusted.}
\end{quote}

In that sense, MASSIF succeeded precisely by disproving its own strongest hypothesis.

\medskip
\hrule
\medskip

\section*{Acknowledgments}

This research was conducted by DUBITO Inc. and Ergo Sum AGI Safety Systems. We acknowledge the contributions of the early CQFT/PFTC research program and thank the anonymous reviewers for their rigorous feedback on earlier drafts.

\medskip

\section*{Conflict of Interest}

The authors declare no competing interests.

\medskip

\section*{Data Availability}

The experimental code and raw data for MASSIF v0.3 are available at the project repository. Trained model weights are from the publicly available GPT-2 model family.

\bibliographystyle{plain}
\bibliography{references}

\medskip
\hrule
\medskip

\begin{thebibliography}{99}

\item[MASSIF v0.3] DUBITO Inc. (2026). \textit{Massive Automated Statistical Sweep Framework} [Computer software]. Version 0.3.

\item[GPT-2] Radford, A., et al. (2019). Language Models are Unsupervised Multitask Learners. \textit{OpenAI Technical Report}.

\item[Transformer Geometry] Vaswani, A., et al. (2017). Attention Is All You Need. \textit{Advances in Neural Information Processing Systems}.

\item[Interpretability] Olah, C., et al. (2020). Zoom In: An Introduction to Circuits. \textit{Distill}.

\item[Latent Space Geometry] Goodfellow, I., et al. (2014). Generative Adversarial Networks. \textit{Advances in Neural Information Processing Systems}.

\end{thebibliography}

\end{document}
